\documentclass[tikz,border=10pt]{standalone}
\usepackage{tikz}
\usepackage{amsmath}
\usepackage{amssymb}
\usepackage{ctex}
\usepackage{pgfplots}
\pgfplotsset{compat=1.18}
\usetikzlibrary{calc, positioning, arrows.meta, decorations.pathreplacing}

\begin{document}
\begin{tikzpicture}[>=Stealth]

%% ===== 图1：原始单位圆 =====
\begin{scope}[xshift=0cm]
  \begin{axis}[
    name=ax1,
    width=5.2cm, height=5.2cm,
    xmin=-1.6, xmax=1.6,
    ymin=-1.6, ymax=1.6,
    axis lines=center,
    xlabel={$x_1$}, ylabel={$x_2$},
    xlabel style={font=\small, right},
    ylabel style={font=\small, above},
    xtick={-1,0,1}, ytick={-1,0,1},
    tick label style={font=\scriptsize},
    grid=both, grid style={gray!15},
    clip=false, axis equal image,
  ]
    % 单位圆
    \addplot[blue!70!black, very thick, fill=blue!10,
             domain=0:360, samples=120]
      ({cos(x)}, {sin(x)});
    % 基向量
    \draw[->,blue!60,thick] (axis cs:0,0) -- (axis cs:1,0)
      node[right,font=\scriptsize] {$\mathbf{e}_1$};
    \draw[->,blue!60,thick] (axis cs:0,0) -- (axis cs:0,1)
      node[above,font=\scriptsize] {$\mathbf{e}_2$};
  \end{axis}
  \node[font=\small\bfseries, blue!70!black,
        fill=blue!10, draw=blue!40, rounded corners=3pt, inner sep=4pt]
    at ([yshift=0.3cm]ax1.north) {原始单位圆};
  \node[font=\scriptsize, gray] at ([yshift=-0.5cm]ax1.south)
    {$\|\mathbf{x}\|=1$};
\end{scope}

%% ===== 箭头1 =====
\node[font=\normalsize\bfseries, orange!80!black]
  at (3.85, 0.05) {$\xrightarrow{V^T}$};

%% ===== 图2：经V^T旋转后（仍是单位圆）=====
\begin{scope}[xshift=4.8cm]
  \begin{axis}[
    name=ax2,
    width=5.2cm, height=5.2cm,
    xmin=-1.6, xmax=1.6,
    ymin=-1.6, ymax=1.6,
    axis lines=center,
    xlabel={$y_1$}, ylabel={$y_2$},
    xlabel style={font=\small, right},
    ylabel style={font=\small, above},
    xtick={-1,0,1}, ytick={-1,0,1},
    tick label style={font=\scriptsize},
    grid=both, grid style={gray!15},
    clip=false, axis equal image,
  ]
    % 单位圆（V^T 是正交矩阵，圆不变形）
    \addplot[green!60!black, very thick, fill=green!10,
             domain=0:360, samples=120]
      ({cos(x)}, {sin(x)});
    % 旋转后的基向量（V^T 旋转 ~30°）
    \draw[->,green!60!black,thick] (axis cs:0,0) -- (axis cs:0.866,0.5)
      node[right,font=\scriptsize] {$\mathbf{v}_1$};
    \draw[->,green!60!black,thick] (axis cs:0,0) -- (axis cs:-0.5,0.866)
      node[above,font=\scriptsize] {$\mathbf{v}_2$};
  \end{axis}
  \node[font=\small\bfseries, green!60!black,
        fill=green!10, draw=green!50, rounded corners=3pt, inner sep=4pt]
    at ([yshift=0.3cm]ax2.north) {$V^T$旋转（仍是圆）};
  \node[font=\scriptsize, gray] at ([yshift=-0.5cm]ax2.south)
    {正交变换保持模长};
\end{scope}

%% ===== 箭头2 =====
\node[font=\normalsize\bfseries, orange!80!black]
  at (8.65, 0.05) {$\xrightarrow{\Sigma}$};

%% ===== 图3：经Σ缩放后椭圆（沿坐标轴）=====
\begin{scope}[xshift=9.6cm]
  \begin{axis}[
    name=ax3,
    width=5.2cm, height=5.2cm,
    xmin=-2.2, xmax=2.2,
    ymin=-1.6, ymax=1.6,
    axis lines=center,
    xlabel={$z_1$}, ylabel={$z_2$},
    xlabel style={font=\small, right},
    ylabel style={font=\small, above},
    xtick={-2,-1,0,1,2}, ytick={-1,0,1},
    tick label style={font=\scriptsize},
    grid=both, grid style={gray!15},
    clip=false, axis equal image,
  ]
    % 椭圆: σ1=1.8, σ2=0.9
    \addplot[orange!80!black, very thick, fill=orange!10,
             domain=0:360, samples=120]
      ({1.8*cos(x)}, {0.9*sin(x)});
    % 半轴标注
    \draw[<->,orange!70,thin]
      (axis cs:0,0) -- (axis cs:1.8,0)
      node[midway,below,font=\scriptsize,orange!80] {$\sigma_1$};
    \draw[<->,orange!70,thin]
      (axis cs:0,0) -- (axis cs:0,0.9)
      node[midway,right,font=\scriptsize,orange!80] {$\sigma_2$};
  \end{axis}
  \node[font=\small\bfseries, orange!80!black,
        fill=orange!10, draw=orange!50, rounded corners=3pt, inner sep=4pt]
    at ([yshift=0.3cm]ax3.north) {$\Sigma$缩放（椭圆，轴对齐）};
  \node[font=\scriptsize, gray] at ([yshift=-0.5cm]ax3.south)
    {奇异值 $\sigma_1\geq\sigma_2>0$};
\end{scope}

%% ===== 箭头3 =====
\node[font=\normalsize\bfseries, orange!80!black]
  at (13.45, 0.05) {$\xrightarrow{U}$};

%% ===== 图4：经U旋转后椭圆（任意方向）=====
\begin{scope}[xshift=14.4cm]
  \begin{axis}[
    name=ax4,
    width=5.2cm, height=5.2cm,
    xmin=-2.2, xmax=2.2,
    ymin=-2.2, ymax=2.2,
    axis lines=center,
    xlabel={$w_1$}, ylabel={$w_2$},
    xlabel style={font=\small, right},
    ylabel style={font=\small, above},
    xtick={-2,-1,0,1,2}, ytick={-2,-1,0,1,2},
    tick label style={font=\scriptsize},
    grid=both, grid style={gray!15},
    clip=false, axis equal image,
  ]
    % 旋转后的椭圆（旋转约40°）
    % 参数: (1.8*cos(t)*cos40 - 0.9*sin(t)*sin40, 1.8*cos(t)*sin40 + 0.9*sin(t)*cos40)
    \addplot[red!70!black, very thick, fill=red!10,
             domain=0:360, samples=120]
      ({1.8*cos(x)*0.766 - 0.9*sin(x)*0.6428},
       {1.8*cos(x)*0.6428 + 0.9*sin(x)*0.766});
    % 主轴方向（u1, u2）
    \draw[->,red!60!black,thick]
      (axis cs:0,0) -- (axis cs:1.38,1.157)
      node[right,font=\scriptsize] {$\mathbf{u}_1$};
    \draw[->,red!60!black,thick]
      (axis cs:0,0) -- (axis cs:-0.579,0.689)
      node[above,font=\scriptsize] {$\mathbf{u}_2$};
  \end{axis}
  \node[font=\small\bfseries, red!70!black,
        fill=red!10, draw=red!40, rounded corners=3pt, inner sep=4pt]
    at ([yshift=0.3cm]ax4.north) {$U$旋转（任意方向椭圆）};
  \node[font=\scriptsize, gray] at ([yshift=-0.5cm]ax4.south)
    {输出 $= A\mathbf{x}$ 的像};
\end{scope}

%% ===== 底部公式框 =====
\node[font=\large, fill=white, draw=gray!50, thin, rounded corners=5pt,
      inner sep=10pt, align=center]
  at (9.6, -2.8) {
    \textbf{奇异值分解（SVD）几何步骤：}
    \quad $A = U\Sigma V^T$\\[4pt]
    $V^T$：旋转输入基\quad
    $\Sigma$：沿坐标轴缩放（奇异值）\quad
    $U$：旋转输出基
  };

%% ===== 整体标题 =====
\node[font=\large\bfseries] at (9.6, 5.0)
  {SVD 三步骤几何：$A = U\Sigma V^T$};

\end{tikzpicture}
\end{document}
